\documentclass[10pt,a4paper]{amsart}
\usepackage[margin=19mm]{geometry}
\usepackage{amsmath,amssymb,mathtools}
\usepackage[scr=rsfs,cal=euler]{mathalfa}
\usepackage{xfrac}
\usepackage[unicode,hidelinks,bookmarks=false]{hyperref}
\newcommand{\ie}{i.e.\ }
\newcommand{\cf}{cf.\ }
\newcommand{\resp}{respectively\ }
\newcommand{\Subsection}{\subsection}
\providecommand{\nobreakdash}{\nobreak\hbox{-}}
\setlength{\emergencystretch}{2em}
\multlinegap0pt
\usepackage{bm}
\usepackage{bbm}
\usepackage{dsfont}
\usepackage{xspace}
\usepackage{cancel}
\usepackage{mathtools}
\usepackage{amsthm}
\makeatletter
\@ifundefined{subsubsection}{\newtheorem{subsubsection}{Subsubsection}}{}
\@ifundefined{cor}{\newtheorem{cor}[subsubsection]{Corollary}}{}
\makeatother
\makeatletter
\newcommand{\Hilbert}{\mathrm{Hilb}}
\newcommand{\Langlands}{\textnormal{\textsf{L}}}
\newcommand{\complexes}{\mathbf C}
\newcommand{\cycle}{\textnormal{\textsf{Z}}}
\newcommand{\deloop}{\textnormal{\textsf{B}}}
\newcommand{\extension}{\textnormal{\textsf{E}}}
\newcommand{\sgn}{\mathrm{sgn}}
\expandafter\ifx\csname void\endcsname\relax
\newenvironment{void}
{\begin{untitledsubsubsection}}
{\end{untitledsubsubsection}}
\fi
\makeatother
\title[What are the extended pure inner forms of a cover?]{What are the extended pure inner forms of a cover?}
\author{Luozi Shi, Yifei Zhao}
\date{Source: arXiv:2506.08696v1 (CC BY 4.0)}
\begin{document}
\maketitle

\noindent\emph{Source and licence.} What are the extended pure inner forms of a cover?. Version arXiv:2506.08696v1 (CC BY 4.0), distributed under the Creative Commons Attribution 4.0 International licence, as recorded in the arXiv licence metadata for this exact version and shown on the abstract page https://arxiv.org/abs/2506.08696v1. Full rights notice: licenses/arxiv-2506.08696-CC-BY-4.0.txt.\par\smallskip

\noindent\textit{Editorial scope.} Counted unit: the Cor whose source label is \texttt{cor-sign-cover-dual-identification} in the article (source0.tex lines 1353--1360, source label \texttt{cor-sign-cover-dual-identification}), delivered together with the article's own complete original proof of it (source0.tex lines 1361--1383, one \texttt{proof} environment of 23 lines). No mathematics is written, repaired or re-derived: statement and proof are the article's own text line for line. Required context retained uncounted: eq-hilbert-cover (lines 890--893); eq-sign-cover-isomorphism-hilbert-cover (lines 1145--1148). Every definition, display and label of those lines is retained unchanged. Omitted from this package and not redistributed: 1--889, 894--1144, 1149--1352, 1384--2752. Nothing inside a retained excerpt is omitted or altered: every retained body is delivered exactly as the article's lines are written, so no omission receipt of any kind accompanies this package. Citations: the retained excerpts carry no citation command at all, so no bibliography entry of the article is transcribed and no reference is redistributed. Reference adaptation: none was needed and none was made; every reference inside the delivered unit resolves inside it, so no unresolved reference or citation is produced. Printed slips kept literal: no printed word, symbol, hypothesis or number of the counted statement or of its proof was altered. Layout and class: the article's own class is \texttt{amsart} with packages including color, hyperref, graphicx, amssymb, epstopdf, enumerate, mathabx, tikz-cd, MnSymbol, mathrsfs; the wrapper is the lane's pinned \texttt{amsart} editorial wrapper at 10pt on A4, which restates the theorem-like environments the retained text needs and adds the article's own macro definitions that the retained text needs (\texttt{\textbackslash Hilbert}, \texttt{\textbackslash Langlands}, \texttt{\textbackslash complexes}, \texttt{\textbackslash cycle}, \texttt{\textbackslash deloop}, \texttt{\textbackslash extension}, \texttt{\textbackslash sgn}), with extra wrapper packages (bm, bbm, dsfont, xspace, cancel, mathtools). The delivered numbering is the unit's own. Assets: no retained span contains a figure environment, a graphics-inclusion command, a PDF-page inclusion, a table or a TikZ picture; no image file is copied into this package, the package declares no asset dependency and no image file is redistributed with it. Licence: version arXiv:2506.08696v1 of this article is distributed under the Creative Commons Attribution 4.0 International licence, as recorded in the arXiv licence metadata for this exact version and shown on the abstract page https://arxiv.org/abs/2506.08696v1. A full rights notice is delivered at licenses/arxiv-2506.08696-CC-BY-4.0.txt. Characters: every retained excerpt is pure ASCII, so the delivered engine, XeLaTeX with its default Latin Modern text font, reports no missing character.
\begin{equation}
\label{eq-hilbert-cover}
1 \rightarrow \{\pm 1\} \rightarrow \widetilde F^{\times}_{\Hilbert} \rightarrow F^{\times} \rightarrow 1
\end{equation}

\begin{equation}
\label{eq-sign-cover-isomorphism-hilbert-cover}
\widetilde F^{\times}_{\sgn} \xrightarrow{\simeq} \widetilde F^{\times}_{\Hilbert}
\end{equation}

\begin{cor}
\label{cor-sign-cover-dual-identification}
There is a canonical isomorphism in $\extension^1(F^{\times}_{/2}, \complexes^{\times})$:
\begin{equation}
\label{eq-sign-cover-dual-identification}
\Langlands_{\deloop\{\pm 1\}}(\widetilde W_F) \xrightarrow{\simeq} \widetilde F^{\times}_{\sgn, /2}.
\end{equation}
\end{cor}

\begin{proof}
For an abelian group $M$, denote by $\cycle^2_e(W_F, M)$ the fiber of the map $e^* : \cycle^2(W_F, M) \rightarrow \cycle^2(*, M)$ given by pullback along the neutral point $e : * \rightarrow */W_F$. Thus $\cycle^2_e(W_F, M)$ is canonically equivalent to the groupoid of central extensions of $W_F$ by $M$.

The restriction of $\Langlands_{\mathbb G_m}$ to $\cycle^2_e(W_F, \complexes^{\times})$ admits the following explicit description: Pulling back a commutative extension of $F^{\times}$ by $\complexes^{\times}$ along the Artin reciprocity map $W_F \rightarrow F^{\times}$ yields an equivalence of groupoids
\begin{equation}
\label{eq-pullback-by-artin-reciprocity}
\extension^1(F^{\times}, \complexes^{\times}) \xrightarrow{\simeq} \cycle^2_e(W_F, \complexes^{\times}),
\end{equation}
whose inverse coincides with the restriction of $\Langlands_{\mathbb G_m}$.

In what follows, we view $\widetilde F^{\times}_{\sgn}$ as an object of $\extension^1(F^{\times}, \complexes^{\times})$. It suffices to identify its image under \eqref{eq-pullback-by-artin-reciprocity} with the extension of $W_F$ by $\complexes^{\times}$ induced from $\widetilde W_F$, and match the $2$-torsion structures defined by $\tau$ and $\widetilde W_F$. The identification follows from the isomorphism \eqref{eq-sign-cover-isomorphism-hilbert-cover}. The matching of $2$-torsion structures follows from an explicit calculation, as we now perform.

Multiplication by $2$ on $\widetilde F^{\times}_{\sgn}$ factors through an isomorphism
\begin{equation}
\label{eq-sgn-pushout-pullback-isomorphism}
\complexes^{\times} \sqcup_{\complexes^{\times}} \widetilde F^{\times}_{\sgn} \xrightarrow{\simeq} \widetilde F^{\times}_{\sgn} \times_{F^{\times}} F^{\times}
\end{equation}
where the push-out is along $(\cdot)^2 : \complexes^{\times} \rightarrow \complexes^{\times}$ and the pullback is along $(\cdot)^2 : F^{\times} \rightarrow F^{\times}$. Using the isomorphism \eqref{eq-sign-cover-isomorphism-hilbert-cover}, we may represent an element of $\widetilde F^{\times}_{\sgn}$ by a pair $(a, z)$ with $a \in F^{\times}$ and $z\in\complexes^{\times}$. Its image under \eqref{eq-sgn-pushout-pullback-isomorphism} is
$$
((a^2, \{a, a\} z^2), a)
$$
which equals the product of $z^2$ with the $(\tau(a), a)$. Hence, splitting of the pushout induced from $\tau$ sends $(a, z)$ to $z^2$. The kernel of this map is the extension \eqref{eq-hilbert-cover}, as desired.
\end{proof}

\end{document}
